% GENERATED FILE. Edit canonical lessons, then run npm run generate:books.
% canonical-source-hash: fnv1a64:71f739b17d77d59e
% canonical-lessons: KA-C419-viphalavagu, KA-C419-phalisu, KA-C419-kaigudu, KA-C419-bikkattu, KA-C419-anivarya

\chapter{Success and Failure}
\label{ch:ka-success-and-failure}
\hlchaptermodality{419}

\begin{chapteropening}
\textbf{By the end of this chapter:} \emph{I can say to fall through, to bear fruit, to come true, a crisis and unavoidable.}

Complete the last lesson of chapter 419: I can say to fall through, to bear fruit, to come true, a crisis and unavoidable.
\end{chapteropening}

\section[\texorpdfstring{\kn{ವಿಫಲವಾಗು} (viphalavāgu)}{ವಿಫಲವಾಗು (viphalavāgu)}]{\kn{ವಿಫಲವಾಗು} (viphalavāgu) — to fall through; to fail}

\label{lesson:KA-C419-viphalavagu}



\begin{quote}
Before the new one: say the Kannada for permanent, then the Kannada for to modify.
\end{quote}

\subsection*{You'll want to know: \kn{ವಿಫಲವಾಗು}}
\textbf{\kn{ವಿಫಲವಾಗು}} — \emph{viphalavāgu} — \textquotedblleft{}to fall through; to fail\textquotedblright{}.

\begin{quote}
\kn{ಮಳೆ} \kn{ಬಂದಿದ್ದರಿಂದ} \kn{ನಮ್ಮ} \kn{ಯೋಜನೆ} \kn{ವಿಫಲವಾಯಿತು}.
\end{quote}

\begin{quote}
\emph{maḷe bandiddarinda namma yōjane viphalavāyitu.}
\end{quote}

\begin{quote}
Because it rained, our plan failed.
\end{quote}

Past (it …): \kn{ವಿಫಲವಾಯಿತು}, \emph{viphalavāyitu}.

\kn{ವಿಫಲ}, fruitless; \kn{ವಿಫಲತೆ}, failure.

\begin{cousinweb}
\textbf{\kn{ವಿಫಲವಾಗು}} (\emph{viphalavāgu}) is built from Sanskrit \emph{viphala}, \textquotedblleft{}fruitless, vain\textquotedblright{} (CDIAL 11803), with native \emph{āgu}, \textquotedblleft{}to become\textquotedblright{} (DEDR 333).
\end{cousinweb}

\begin{grammarlens}[title={What you've built}]
One more word for saying what happens if.
\end{grammarlens}

\subsection*{Your turn}
\begin{itemize}
\raggedright
  \item \emph{Say it:} \emph{viphalavāgu}
  \item \emph{Say it:} \emph{viphalavāgu}, once more
  \item \emph{Say it:} \emph{maḷe bandiddarinda namma yōjane viphalavāyitu.} — Because it rained, our plan failed.
  \item \emph{Recall:} say the Kannada for permanent, then the Kannada for to modify, then say \emph{viphalavāgu} again
\end{itemize}

\begin{tcolorbox}[breakable,colback=teal!4,colframe=teal!35!black,title={Before you move on}]
What does \kn{ವಿಫಲವಾಗು} mean? (\textquotedblleft{}To fall through; to fail\textquotedblright{}.) Say it once more.
\end{tcolorbox}
\section[\texorpdfstring{\kn{ಫಲಿಸು} (phalisu)}{ಫಲಿಸು (phalisu)}]{\kn{ಫಲಿಸು} (phalisu) — to bear fruit}

\label{lesson:KA-C419-phalisu}



\begin{quote}
Before the new one: say the Kannada for to modify, then the Kannada for to fall through.
\end{quote}

\subsection*{You'll want to know: \kn{ಫಲಿಸು}}
\textbf{\kn{ಫಲಿಸು}} — \emph{phalisu} — \textquotedblleft{}to bear fruit\textquotedblright{}.

\begin{quote}
\kn{ಅವಳ} \kn{ಪ್ರಯತ್ನ} \kn{ಕೊನೆಗೆ} \kn{ಫಲಿಸಿತು}.
\end{quote}

\begin{quote}
\emph{avaḷa prayatna konege phalisitu.}
\end{quote}

\begin{quote}
Her effort finally paid off.
\end{quote}

Past (it …): \kn{ಫಲಿಸಿತು}, \emph{phalisitu}.

Of efforts, prayers and trees.

\begin{cousinweb}
\textbf{\kn{ಫಲಿಸು}} (\emph{phalisu}) is built on Sanskrit \emph{phala}, \textquotedblleft{}fruit\textquotedblright{} (CDIAL 9051), with the Kannada verb ending -\emph{isu}.
\end{cousinweb}

\begin{grammarlens}[title={What you've built}]
One more word for saying what happens if.
\end{grammarlens}

\subsection*{Your turn}
\begin{itemize}
\raggedright
  \item \emph{Say it:} \emph{phalisu}
  \item \emph{Say it:} \emph{phalisu}, once more
  \item \emph{Say it:} \emph{avaḷa prayatna konege phalisitu.} — Her effort finally paid off.
  \item \emph{Recall:} say the Kannada for to modify, then the Kannada for to fall through, then say \emph{phalisu} again
\end{itemize}

\begin{tcolorbox}[breakable,colback=teal!4,colframe=teal!35!black,title={Before you move on}]
What does \kn{ಫಲಿಸು} mean? (\textquotedblleft{}To bear fruit\textquotedblright{}.) Say it once more.
\end{tcolorbox}
\section[\texorpdfstring{\kn{ಕೈಗೂಡು} (kaigūḍu)}{ಕೈಗೂಡು (kaigūḍu)}]{\kn{ಕೈಗೂಡು} (kaigūḍu) — to come true; to work out}

\label{lesson:KA-C419-kaigudu}



\begin{quote}
Before the new one: say the Kannada for to fall through, then the Kannada for to bear fruit.
\end{quote}

\subsection*{You'll want to know: \kn{ಕೈಗೂಡು}}
\textbf{\kn{ಕೈಗೂಡು}} — \emph{kaigūḍu} — \textquotedblleft{}to come true; to work out\textquotedblright{}.

\begin{quote}
\kn{ಮನೆ} \kn{ಕಟ್ಟುವ} \kn{ಅವರ} \kn{ಕನಸು} \kn{ಕೈಗೂಡಿತು}.
\end{quote}

\begin{quote}
\emph{mane kaṭṭuva avara kanasu kaigūḍitu.}
\end{quote}

\begin{quote}
Their dream of building a house came true.
\end{quote}

Past (it …): \kn{ಕೈಗೂಡಿತು}, \emph{kaigūḍitu}.

Of plans and wishes, never of people.

\begin{cousinweb}
\textbf{\kn{ಕೈಗೂಡು}} (\emph{kaigūḍu}) is native Dravidian, built from \emph{kai}, \textquotedblleft{}hand\textquotedblright{} (DEDR 2023), and \emph{kūḍu}, \textquotedblleft{}to join\textquotedblright{} (DEDR 1882): to come into the hand.
\end{cousinweb}

\begin{grammarlens}[title={What you've built}]
One more word for saying what happens if.
\end{grammarlens}

\subsection*{Your turn}
\begin{itemize}
\raggedright
  \item \emph{Say it:} \emph{kaigūḍu}
  \item \emph{Say it:} \emph{kaigūḍu}, once more
  \item \emph{Say it:} \emph{mane kaṭṭuva avara kanasu kaigūḍitu.} — Their dream of building a house came true.
  \item \emph{Recall:} say the Kannada for to fall through, then the Kannada for to bear fruit, then say \emph{kaigūḍu} again
\end{itemize}

\begin{tcolorbox}[breakable,colback=teal!4,colframe=teal!35!black,title={Before you move on}]
What does \kn{ಕೈಗೂಡು} mean? (\textquotedblleft{}To come true; to work out\textquotedblright{}.) Say it once more.
\end{tcolorbox}
\section[\texorpdfstring{\kn{ಬಿಕ್ಕಟ್ಟು} (bikkaṭṭu)}{ಬಿಕ್ಕಟ್ಟು (bikkaṭṭu)}]{\kn{ಬಿಕ್ಕಟ್ಟು} (bikkaṭṭu) — a crisis; a deadlock}

\label{lesson:KA-C419-bikkattu}



\begin{quote}
Before the new one: say the Kannada for to bear fruit, then the Kannada for to come true.
\end{quote}

\subsection*{You'll want to know: \kn{ಬಿಕ್ಕಟ್ಟು}}
\textbf{\kn{ಬಿಕ್ಕಟ್ಟು}} — \emph{bikkaṭṭu} — \textquotedblleft{}a crisis; a deadlock\textquotedblright{}.

\begin{quote}
\kn{ನೀರಿನ} \kn{ಬಿಕ್ಕಟ್ಟಿನಿಂದ} \kn{ಊರಿನ} \kn{ಜನರು} \kn{ಕಷ್ಟಪಟ್ಟರು}.
\end{quote}

\begin{quote}
\emph{nīrina bikkaṭṭininda ūrina janaru kaṣṭapaṭṭaru.}
\end{quote}

\begin{quote}
The townspeople suffered from the water crisis.
\end{quote}

\kn{ಬಿಕ್ಕಟ್ಟಿಗೆ} \kn{ಸಿಕ್ಕು}, to land in a fix.

\begin{cousinweb}
\textbf{\kn{ಬಿಕ್ಕಟ್ಟು}} (\emph{bikkaṭṭu}) is a Kannada word whose origin is not settled: DEDR does not list it.
\end{cousinweb}

\begin{grammarlens}[title={What you've built}]
One more word for saying what happens if.
\end{grammarlens}

\subsection*{Your turn}
\begin{itemize}
\raggedright
  \item \emph{Say it:} \emph{bikkaṭṭu}
  \item \emph{Say it:} \emph{bikkaṭṭu}, once more
  \item \emph{Say it:} \emph{nīrina bikkaṭṭininda ūrina janaru kaṣṭapaṭṭaru.} — The townspeople suffered from the water crisis.
  \item \emph{Recall:} say the Kannada for to bear fruit, then the Kannada for to come true, then say \emph{bikkaṭṭu} again
\end{itemize}

\begin{tcolorbox}[breakable,colback=teal!4,colframe=teal!35!black,title={Before you move on}]
What does \kn{ಬಿಕ್ಕಟ್ಟು} mean? (\textquotedblleft{}A crisis; a deadlock\textquotedblright{}.) Say it once more.
\end{tcolorbox}
\section[\texorpdfstring{\kn{ಅನಿವಾರ್ಯ} (anivārya)}{ಅನಿವಾರ್ಯ (anivārya)}]{\kn{ಅನಿವಾರ್ಯ} (anivārya) — unavoidable}

\label{lesson:KA-C419-anivarya}



\begin{quote}
Before the new one: say the Kannada for to come true, then the Kannada for a crisis.
\end{quote}

\subsection*{You'll want to know: \kn{ಅನಿವಾರ್ಯ}}
\textbf{\kn{ಅನಿವಾರ್ಯ}} — \emph{anivārya} — \textquotedblleft{}unavoidable\textquotedblright{}.

\begin{quote}
\kn{ಮಳೆ} \kn{ಬಂದರೆ} \kn{ಪಂದ್ಯ} \kn{ಮುಂದೂಡುವುದು} \kn{ಅನಿವಾರ್ಯ}.
\end{quote}

\begin{quote}
\emph{maḷe bandare pandya mundūḍuvudu anivārya.}
\end{quote}

\begin{quote}
If it rains, putting off the match is unavoidable.
\end{quote}

\kn{ಅನಿವಾರ್ಯವಾಗಿ}, of necessity.

\begin{cousinweb}
\textbf{\kn{ಅನಿವಾರ್ಯ}} (\emph{anivārya}) is borrowed from Sanskrit \emph{anivārya}, \textquotedblleft{}not to be warded off, unavoidable\textquotedblright{} (Monier-Williams).
\end{cousinweb}

\begin{grammarlens}[title={What you've built}]
One more word for saying what happens if.
\end{grammarlens}

\subsection*{Your turn}
\begin{itemize}
\raggedright
  \item \emph{Say it:} \emph{anivārya}
  \item \emph{Say it:} \emph{anivārya}, once more
  \item \emph{Say it:} \emph{maḷe bandare pandya mundūḍuvudu anivārya.} — If it rains, putting off the match is unavoidable.
  \item \emph{Recall:} say the Kannada for to come true, then the Kannada for a crisis, then say \emph{anivārya} again
\end{itemize}

\begin{tcolorbox}[breakable,colback=teal!4,colframe=teal!35!black,title={Before you move on}]
What does \kn{ಅನಿವಾರ್ಯ} mean? (\textquotedblleft{}Unavoidable\textquotedblright{}.) Say it once more.
\end{tcolorbox}
